\documentclass[11pt]{article}

\usepackage[a4paper,margin=25mm]{geometry}
\usepackage{amsmath,amssymb}
\usepackage{booktabs}
\usepackage{enumitem}
\usepackage{tikz}
\usepackage[colorlinks=true,linkcolor=blue,citecolor=blue,urlcolor=blue]{hyperref}

\title{\textbf{SCT--CG1 Theory Manual: Crack Case TC23B}\\[4pt]
\large Cracked Stringer at a Fastener Hole with Skin Bridging}
\author{Southward Consulting Ltd}
\date{June 2026 \quad(rev.\ A)}

% lightweight "Assumptions" list
\newenvironment{assume}{\par\smallskip\noindent\textit{Assumptions:}\begin{itemize}[leftmargin=1.5em,topsep=2pt,itemsep=1pt]}{\end{itemize}\par\smallskip}

\newcommand{\ksi}{\,\mathrm{ksi}}
\newcommand{\inch}{\,\mathrm{in}}

\begin{document}
\maketitle

\begin{abstract}
TC23B extends NASGRO crack case TC23 (unequal through cracks at an offset hole
in a finite-width plate) to a stringer, modelled by its unfolded width, that is
attached to a background fuselage skin by a single vertical row of fasteners
through the critical hole. As the stringer crack opens, the fasteners shed load
into the intact skin, which bridges the crack and retards growth. This manual
documents the complete theory and, in dedicated sections, every assumption made
in the application of (i) bending restraint, (ii) the bearing-load model
options, and (iii) the skin far-field stress state. It also documents the
skin-bridging formulation, its two stress-intensity-reduction routes, and the
FEA-derived correction surface used in the post-link-up phase. LEFM and the
NASGRO crack-growth equation are assumed known and are summarised only briefly.
\end{abstract}

\tableofcontents

% ====================================================================
\section{Scope, conventions and nomenclature}

\subsection{Configuration}
The cracked member is a stringer represented by its \emph{unfolded width} $W$
(a flat membrane of thickness $t$). A fastener hole of diameter $D$ lies a
distance $m$ from one edge (the right margin). Two through cracks $c_1$ (wide
side) and $c_2$ (edge side) grow from the hole. A single vertical row of
fasteners at pitch $p$ passes through the hole, parallel to the load, attaching
the stringer to a background \emph{infinite} skin (modulus $E_2$, thickness
$t_2$). Loading comprises remote tension $S_0$, in-plane bending $S_2$, and
pin/bearing $S_3$.

\subsection{Conventions}
\begin{itemize}[leftmargin=1.5em,topsep=2pt,itemsep=1pt]
\item Units are consistent Imperial: ksi, in, kip (and ksi$\sqrt{\text{in}}$ for $K$).
\item Hole offset: $m$ = hole-centre to right edge; $B = W-m$ = hole-centre to
  left edge; eccentricity $e = W/2 - B = m - W/2$.
\item Plane stress throughout. Stress ratio $R = \sigma_{\min}/\sigma_{\max}$.
\item $S_2>0$ places tension at the crack mouth (opens the crack); $S_3$ is the
  bearing stress $P/(Dt)$, clipped to $\ge 0$.
\end{itemize}

% ====================================================================
\section{Baseline stress intensity (inherited TC23)}

For each crack tip $i$ the SIF is the NASGRO Appendix~C compounded solution
\begin{equation}
K_i = \Bigl(\beta_i^{A} S_0 + \beta_i^{B} S_2 + b_{p,i}\,\bar S_3\Bigr)\sqrt{\pi c_i},
\label{eq:Ktc23}
\end{equation}
with the remote-tension factor $\beta^A = \beta^{A1}\beta^{A2}\beta^{A3}$
(Tweed--Rooke/Bowie unequal-crack fit $\times$ the $\varphi_1/\varphi_2$
finite-width corrections), the in-plane-bending factor $\beta^B$ (Solution~B,
\S\ref{sec:bend}), and the bearing factor $b_p$ (Solution~C, \S\ref{sec:bearing}).
Equation~\eqref{eq:Ktc23} is verified term-by-term against the NASGRO~8.2
manual; this document treats it as the validated baseline onto which the TC23B
extensions are applied.

% ====================================================================
\section{Bending restraint}
\label{sec:bend}

Two physically distinct ``bending'' effects appear in TC23B and must not be
conflated: (a) \emph{applied in-plane bending} $S_2$ of the stringer
cross-section, and (b) \emph{load-line-eccentricity bending} of the post-link-up
single-edge-notch-tension (SENT) ligament, which is what the bending-restraint
parameter $\eta$ controls.

\subsection{Applied in-plane bending, $S_2$ (Solution B)}
In the two-crack phase the bending factor uses NASGRO TC23 Solution~B:
$\beta_i^{B}=\beta_i^{B1}\,\beta_i^{B2}$, where $\beta^{B1}$ is the Bowie hole
factor (=$\beta^{A1}$) and $\beta^{B2}=\Phi_1(d_1,d_2)$ is interpolated from the
NASGRO FEM benchmark table for a crack of size $2c_0$ in a finite plate under
in-plane bending, with $d_1=(b-c_0)/W$, $d_2=(b+c_0)/W$ and
$\beta_2^{B2}=-\Phi_1(1-d_2,1-d_1)$. In the post-link-up phase $S_2$ uses the
TC02 Solution~B (Gross--Srawley) edge-crack factor
\begin{equation}
F_B(\lambda)=1.122-1.40\lambda+7.33\lambda^2-13.08\lambda^3+14.0\lambda^4,
\qquad \lambda=a/W .
\end{equation}

\begin{assume}
\item $S_2$ is the \emph{extreme-fibre} in-plane bending stress at the cracked
  cross-section ($=6M/BW^2$); the linear stress distribution is integrated
  internally by $\Phi_1$ / $F_B$, so only the surface value is supplied.
\item Sign: $S_2>0$ opens the crack and is added with its sign for any
  eccentricity $e$. The TC23 two-crack convention (tension on the right face)
  is continuous with the post-link-up SENT crack mouth at the right edge.
\item $\Phi_1$ is bilinearly interpolated from the discrete FEM grid; values
  outside the grid are clamped to the grid edge.
\item In-plane bending is treated as an independent load system superposed
  linearly on tension and bearing (LEFM superposition); no plasticity or
  load-redistribution coupling between systems is modelled.
\end{assume}

\subsection{Load-line-eccentricity bending of the SENT ligament ($\eta$)}
After link-up the remaining ligament is an edge crack growing across the
stringer. Because the crack offsets the load path, an SENT coupon develops
secondary bending; the resulting tip SIF depends on how strongly the
surrounding structure restrains the rotation. This is captured by linearly
interpolating between the two limiting Tada/Paris/Irwin SENT solutions:
\begin{align}
F_{\text{unrestr}}(\lambda) &= 1.122-0.231\lambda+10.550\lambda^2-21.710\lambda^3+30.382\lambda^4 \quad(\eta=0),\\
F_{\text{restr}}(\lambda)   &= 1.122-0.561\lambda-0.205\lambda^2+0.471\lambda^3-0.190\lambda^4 \quad(\eta=1),\\
\beta_{\text{SENT}}(\lambda,\eta) &= F_{\text{unrestr}} + \eta\,(F_{\text{restr}}-F_{\text{unrestr}}).
\label{eq:bsent}
\end{align}

\begin{assume}
\item $\eta\in[0,1]$ is a user input. $\eta=0$ = pinned/free ends (free
  rotation, maximum secondary bending, highest $K$); $\eta=1$ =
  clamped/built-in ends (rotation suppressed, lowest $K$). Since
  $F_{\text{restr}}<F_{\text{unrestr}}$, increasing $\eta$ monotonically reduces
  $\beta_{\text{SENT}}$ and $K$.
\item The interpolation is \emph{linear} in $\eta$; intermediate $\eta$ has no
  independent solution and is a pragmatic blend, not an exact boundary-value
  result. The user must choose $\eta$ to represent the actual rotational
  restraint of the installation (e.g.\ continuity of the stringer beyond the
  bay, frame/clip attachment).
\item The two end-member polynomials carry their handbook accuracy ($\approx
  0.5\%$; the unrestrained fit is quoted to $a/W\le0.6$, the restrained fit for
  any $a/W$). Accuracy of intermediate $\eta$ is not separately validated.
\item $\eta$ acts only in the post-link-up SENT phase. In the two-crack phase
  the symmetric hole geometry carries no load-line eccentricity and $\eta$ is
  not applied.
\item $\eta$ multiplies the membrane SENT factor only; the separate $S_2$
  in-plane-bending term ($F_B$) is scaled by the same stiffener-shielding factor
  but is not a function of $\eta$.
\item Default $\eta=0$ (conservative: maximum SENT bending).
\end{assume}

% ====================================================================
\section{Bearing-load model options}
\label{sec:bearing}

The bearing (pin-load) factor $b_p$ such that $K_{\text{bear}}=b_p\,\bar
S_3\sqrt{\pi c}$ is selectable. The effective pin-load stress is clipped,
$\bar S_3=\max(S_3,0)$, and the gross-section stress used elsewhere is
$S_{\text{gross}}=S_0+(D/W)\bar S_3$.

\subsection{Option 1 --- NASGRO TC23 Solution C (default)}
The Appendix~C compounding $b_p=(D/W)\,\beta^C$, with
$\beta_i^C=\tfrac12(\beta_i^{C1}+\beta_i^{C2}+\beta_i^{C3})$; $\beta^{C1}=\beta^A$;
$\beta^{C2}=\beta^{C2{:}1}\beta^{C2{:}2}\beta^{C2{:}3}\beta^{C2{:}4}$ with the
radial-pressure factor $e^{0.15(\rho_i^2-1)}$, the unequal-crack factor
$\beta^{C2{:}3}=\tfrac{W}{\pi c_0}\sqrt{(c_0\mp\zeta)/(c_0\pm\zeta)}$,
$\zeta=(c_1-c_2)/2$, and the finite-width factor $\beta^{C2{:}4}$; and
$\beta^{C3}=(6e/W)\beta^{B1}\beta^{B2}$ (the bending raised by the pin
eccentricity).
\begin{assume}
\item The pin load is idealised as a point wedge force on the crack faces (the
  $W/\pi c_0$ term). This is asymptotically correct for cracks long relative to
  the hole ($c\gtrsim 2D$) but \emph{conservative at short cracks}, where the
  real distributed bore-pressure field decays faster than a point force.
\item Reacted on one side (one-sided reaction), consistent with NASGRO; the
  bearing contribution adds $(D/W)\bar S_3$ to the gross stress.
\item Compression is clipped ($\bar S_3\ge0$): a pin cannot pull on the bore.
\end{assume}

\subsection{Option 2 --- Hybrid TC05 FEM $\rightarrow$ TC23}
NASGRO TC05 FEM pin-load beta tables (Table~C7, two cracks at a loaded hole,
strip$\leftrightarrow$row equivalence $D/H=D/W$, $x=c/(W-D)$) are used at short
cracks, blended linearly in $c/D$ to Solution~C over the band $c/D=0.5\ldots1.0$:
\begin{equation}
b_p(c)=(1-T)\,F_{3,\text{TC05}}\!\Bigl(\tfrac{c}{W-D},\tfrac{D}{W}\Bigr)
       + T\,(D/W)\beta^C_{\text{TC23}},\qquad
T=\mathrm{clamp}\!\Bigl(\tfrac{c/D-0.5}{0.5},0,1\Bigr).
\end{equation}
\begin{assume}
\item The TC05 FEM data are the more accurate short-crack solution (they capture
  the bore-field decay); the point-force Solution~C is more accurate once the
  crack is long relative to the bore and carries the offset/finite-width effects
  TC05 lacks. The blend band ($0.5D$--$1.0D$) is a modelling choice, not derived.
\item The strip is mapped to the TC05 hole row at pitch $H=W$; $D/H$ and the
  normalised crack length are clamped to the table ranges.
\item Continuity of $b_p(c)$ through the band is enforced and regression-tested.
\end{assume}

\subsection{Option 3 --- Effective width ($W_{\text{eff}}=7D$)}
The transferred load is reacted locally over an effective width of a few
diameters (the $\sim35^\circ$ Inglis load-spread), so the bearing case behaves
like a remote tension acting on the cracked hole:
\begin{equation}
b_p(c)=\frac{D}{W_{\text{eff}}}\,\beta^A(c),\qquad
W_{\text{eff}}=\min(k_{\text{brg}}D,\,W),\quad k_{\text{brg}}=7 .
\end{equation}
\begin{assume}
\item The bearing load does not see plate edges farther than $\sim$$W_{\text{eff}}/2$
  (AFGROW wide-panel guidance: keep the hole within $\sim6D$ of an edge for
  bearing). $k_{\text{brg}}=7$ is the default and is user-overridable.
\item The bearing beta is obtained by scaling the full-width tension beta
  $\beta^A$ by $D/W_{\text{eff}}$; it therefore inherits $\beta^A$'s
  long-crack edge magnification (a known difference from a beta computed in the
  $W_{\text{eff}}$ sub-geometry).
\item This option reproduces an external benchmark's bearing betas to within a
  few percent and is the least conservative of the three at short cracks.
\end{assume}

\paragraph{Ordering.} At the geometries examined the three options rank
NASGRO Solution~C (most severe) $>$ Hybrid $>$ Effective-width (least severe) at
short cracks; they cross over at long cracks. The choice is the analyst's and
should be recorded in the substantiation.

% ====================================================================
\section{Skin far-field stress state}
\label{sec:skin}

The skin's longitudinal strain relative to the stringer governs how much load
the fasteners transfer in the far field. Three modes are provided, entering the
bridging compatibility solve (\S\ref{sec:bridge}) as a relative slip
$\Delta\varepsilon\,y_i$ per fastener (measured from the crack plane). Under
\emph{proportional loading} the skin stresses are specified at the peak stringer
stress $\hat\sigma$ and scale with it through the cycle.

\begin{description}[leftmargin=1.5em]
\item[Strain-matched (default).] $\Delta\varepsilon=0$. The skin and stringer
  far-field longitudinal strains are equal (edges effectively coincident);
  fastener loads arise from crack opening only.
\item[Strain-matched $+$ hoop $\sigma_H$ (Poisson).] Longitudinal compatibility
  is retained and a transverse (hoop) stress contracts the skin via Poisson:
  \begin{equation}
  \Delta\varepsilon/\sigma=\frac{\nu_{\text{skin}}\,\sigma_H}{E_{\text{skin}}\,\hat\sigma}.
  \end{equation}
\item[Fully specified $(\sigma_L,\sigma_H)$.]
  \begin{equation}
  \Delta\varepsilon/\sigma=\frac{1}{E_{\text{str}}}
   -\frac{\sigma_L-\nu_{\text{skin}}\sigma_H}{E_{\text{skin}}\,\hat\sigma}.
  \end{equation}
  Equivalent to the hoop mode when $\sigma_L=\hat\sigma\,E_{\text{skin}}/E_{\text{str}}$.
\end{description}

\begin{assume}
\item \textbf{Proportional loading}: $\sigma_L,\sigma_H$ scale with the stringer
  stress, so the strain mismatch per unit stress is constant; the bridging
  restraint is then stress-independent and the constant-amplitude integration is
  unchanged. Non-proportional skin loading is not modelled.
\item \textbf{Input stress is the stringer far-field stress.} $\hat\sigma$ is the
  stress actually in the stringer remote from the joint, already net of any
  global stiffness-ratio load sharing. Consequently the strain mismatch enters
  bridging only through the crack-\emph{driven} fastener-load increment: the
  no-crack baseline $\sum_j F_j(0)$ is subtracted (\S\ref{sec:bridge}).
\item \textbf{Second-order effect on life.} Because its primary action is a
  uniform far-field load share already present in $\hat\sigma$, the biaxial/hoop
  state has only a few-percent effect on crack-growth life through the bridging
  route. To make a hoop/biaxial load share retard the crack as a first-order
  effect it must instead be applied as a reduction of the input far-field
  stringer stress (a separate stiffness-ratio calculation outside this tool).
\item ``Fully specified'' with $\sigma_L=0$ represents a longitudinally
  \emph{unloaded} skin (a pure doubler) --- a valid but distinct case from the
  hoop case; the warning in the input reflects this.
\item Load attraction (skin strained more than the stringer, large $\sigma_L$)
  reverses the mismatch sign and is admitted (the restraint factor may exceed
  unity, clamped to $[0.05,2]$).
\item Plane-stress, linear-elastic skin; in-plane (not out-of-plane / pillowing)
  behaviour only.
\end{assume}

% ====================================================================
\section{Skin bridging}
\label{sec:bridge}

\subsection{Displacement-compatibility solve}
By symmetry the fasteners act in pairs at $y_j=\pm jp$ (the fastener in the
cracked hole lies on the crack plane and carries no bridging load). With pair
forces $F_j$ and fastener flexibility $f$, compatibility of the relative
stringer--skin displacement at each fastener gives
\begin{equation}
\bigl[f\,\mathbf I + \mathbf G^{\text{str}} + \mathbf G^{\text{sk}}\bigr]\mathbf F
   = \mathbf V_0\,\sigma + \Delta\varepsilon\,\mathbf y,
\label{eq:compat}
\end{equation}
solved each step at the current effective half-crack $a$. $V_0$ is the
Westergaard crack-opening of the stringer at each fastener; $\mathbf G^{\text{sk}}$
is the intact-skin 2-D Kelvin (Muskhelishvili) point-force compliance;
$\mathbf G^{\text{str}}$ is the cracked-sheet compliance (uncracked Kelvin field
plus the exact crack-release correction from the Betti/Rice weight-function
identity $v^{\text{corr}}_{ij}=-(t_1/E_1)\int_0^a[K^+_iK^+_j+K^-_iK^-_j]\,da'$).

\subsection{Fastener flexibility --- modified Tate \& Rosenfeld}
\begin{equation}
f=\frac{0.375}{E_f t_1}+\frac{0.375}{E_f t_2}
 +\frac{0.9}{E_1 t_1}+\frac{0.9}{E_2 t_2}
 +\frac{32(1+\nu_f)(t_1+t_2)}{9E_f\pi d^2}
 +\frac{8\,(t_1^3+5t_1^2t_2+5t_1t_2^2+t_2^3)}{5E_f\pi d^4}\;\;[\text{in/kip}].
\end{equation}

\subsection{Reducing the tip SIF --- two routes}
The solved fastener loads reduce the membrane (tension) contribution to $K$.
Two routes are implemented; they account for the \emph{same} $F_j$ and must not
both be applied.
\begin{description}[leftmargin=1.5em]
\item[Weight-function (crack-tip shielding).]
  $R=\bigl(\sqrt{\pi a}+\sum_j F_j K^+_j(a)\bigr)/\sqrt{\pi a}$. Each fastener's
  contribution is weighted by its centre-crack SIF Green's function $K^+_j$, so
  near-tip fasteners shield more than far ones.
\item[Load bypass (load-path).]
  $R_{\text{bp}}=1-P_{\text{tot}}/(S_{\text{gross}}Wt)$ with
  $P_{\text{tot}}=\sum_j F_j(a)-\sum_j F_j(0)$ (the crack-driven one-sided
  bypassed load; the no-crack baseline is subtracted). The membrane $K$ uses the
  bypass stress $S_{\text{byp}}=R_{\text{bp}}S_0$; bending and bearing keep $S_0$.
\end{description}
Both are applied to the membrane term only:
$K=R\,\beta\,S_0\sqrt{\pi c}+K_{\text{bend}}+K_{\text{bear}}$, with $\beta$ the
pure geometry factor (the restraint shows as a separate $S_{\text{byp}}/S_0$
column in the log). $R$ is clamped to $[0.05,2]$.

\begin{assume}
\item The displacement-compatibility solve uses infinite-sheet elastic fields
  for the stringer; finite width enters only through the TC23 $\beta$ factors
  (standard compounding). Plane stress throughout.
\item The hole is omitted from the bridging displacement field (centre-crack
  idealisation $2c_0=c_1+D+c_2$); the fastener line passes through the flaw
  centre (exact for $c_1=c_2$).
\item The Kelvin self-term is regularised over a cut-off radius $d/2$;
  deformation inside that radius is represented by the empirical fastener
  flexibility, not double-counted.
\item Fastener loads scale linearly with stress $\Rightarrow R$ is
  stress-independent (constant-amplitude integration unaffected).
\item Fasteners are XY load-transfer springs (single shear flexibility $f$); no
  preload, friction, bearing nonlinearity, or progressive fastener failure
  (an optional shear-allowable check only flags exceedance).
\end{assume}

\subsection{Validation and the FEA correction surface}
\label{sec:fea}
An independent finite-element model (FEA-membranes, interaction-integral $K$,
validated to $-0.1\%$ on SENT/CCT handbook benchmarks) of the cracked-stringer +
skin + fastener-row configuration gives the true restraint ratio
$R_{\text{FEA}}=K_{\text{with fasteners}}/K_{\text{without fasteners}}$. At the
anchor point ($a/W=0.35$) it gives $R_{\text{FEA}}=0.811$, against which the
weight-function route ($0.819$, $+1.0\%$) is far more accurate than the load-bypass
route ($0.886$, $+9.3\%$). The weight-function (crack-tip-shielding) route is
therefore the validated analytical method; the load-bypass route under-predicts
the bridging (higher $K$, shorter life) in the SENT phase.

For the post-link-up SENT phase a 1-D FEA \emph{correction surface}
$R(a/W)=K_{\text{bridged}}/K_{\text{unbridged}}$ is tabulated for a fixed
configuration and interpolated at run time (the same FEM-table approach used for
$\Phi_1$ bending and TC05 bearing): $K_{\text{SENT}}=R(a/W)\,\beta_{\text{SENT}}
S_0\sqrt{\pi a}$.
\begin{assume}
\item The FEA restraint ratio cancels mesh and SENT-$\beta$ idealisation (a
  ratio from the same mesh), so it isolates the bridging knock-down.
\item The 1-D table is valid only for its tabulated configuration (geometry,
  fastener stiffness, skin/stringer stiffness ratio); a geometry mismatch is
  flagged. Outside the tabulated $a/W$ range the endpoints are held (flat
  extrapolation).
\item The benchmark FEA panel uses a near-uniform crack-plane field (St-Venant);
  the dimensionless restraint ratio is insensitive to the exact far-field load
  introduction.
\end{assume}

% ====================================================================
\section{Post-link-up SENT phase}
When the edge-side ligament fails ($c_2$ reaches the plate edge) the analysis
continues as a single edge crack of depth $a_{\text{edge}}=c_1+D+(W/2-e-R)$
growing from the right edge. The tip SIF combines tension and bearing on the
membrane SENT factor $\beta_{\text{SENT}}(a/W,\eta)$ and bending on $F_B$,
reduced by the bridging restraint (\S\ref{sec:bridge}, or the FEA surface
\S\ref{sec:fea}). In the bridging model the edge crack is idealised as a half
centre-crack centred on the plate edge, fastener line at depth $m$, with
free-edge image sources at $-m$.

% ====================================================================
\section{Crack-growth integration}
The NASGRO equation
$da/dN=C[(1-f_N)/(1-R)\,\Delta K]^n(1-\Delta K_{th}/\Delta K)^p/(1-K_{\max}/K_c)^q$
is integrated with adaptive crack-step control; both tips advance independently
with a shared $\Delta N$. Plastic-zone correction has three modes:
\emph{off}; \emph{ligament} (default --- the Irwin correction activates per tip
only once the ligament net-section stress $\sigma\,\text{trib}/(\text{trib}-R-c)$
reaches yield); \emph{always}. Termination: fracture ($K_{\max}\ge K_c$),
net-section yield on flow stress, geometry limit, or max cycles. Results are
logged at a user cycle interval.

\begin{assume}
\item Constant-amplitude loading; no load-interaction/retardation beyond the
  selected PZC mode.
\item Net-section yield uses the gross stress (including the bearing-equivalent
  membrane term) and ignores the load shed to the skin (conservative).
\end{assume}

% ====================================================================
\section{Consolidated assumptions and limitations}
\begin{enumerate}[leftmargin=1.6em,topsep=2pt,itemsep=1pt]
\item Linear-elastic plane stress; isotropic sheets; small-displacement LEFM
  with linear superposition of $S_0$, $S_2$, $S_3$ and bridging.
\item TC23 compounded $\beta$ factors carry their NASGRO Appendix~C accuracy;
  finite width enters bridging only through them.
\item Bending restraint $\eta$ is a linear blend of two SENT end-members and a
  user judgement of installation rotational restraint (\S\ref{sec:bend}).
\item Bearing model is analyst-selected among three options with the documented
  short/long-crack trade-offs (\S\ref{sec:bearing}).
\item Skin far-field state assumes proportional loading and that the input
  stringer stress is already net of global load sharing (\S\ref{sec:skin}).
\item The validated bridging $K$-reduction is the weight-function route; the
  SENT phase may instead use the FEA correction surface (\S\ref{sec:fea}).
\item No progressive fastener failure, contact/friction, residual stress,
  out-of-plane (pillowing) bending, or environmental effects.
\end{enumerate}

% ====================================================================
\begin{thebibliography}{9}
\bibitem{nasgro} NASGRO Reference Manual v8.2, SwRI --- Appendix~C, crack cases
  TC23, TC05, TC02.
\bibitem{tada} H.~Tada, P.~C. Paris, G.~R. Irwin, \emph{The Stress Analysis of
  Cracks Handbook}, 3rd ed., ASME Press, 2000 (SENT and edge-crack factors).
\bibitem{gross} B.~Gross, J.~E. Srawley, NASA TN~D-2603, 1965 (single-edge-notch
  bending/tension factors).
\bibitem{swift} T.~Swift, ``Fracture Analysis of Stiffened Structure,'' ASTM
  STP~842, 1984; ``Repairs to Damage Tolerant Aircraft,'' FAA-AIR-90-01, 1990.
\bibitem{tate} M.~B. Tate, S.~J. Rosenfeld, NACA TN-1051, 1946 (fastener
  flexibility; modified form with 0.375/0.9 hole-compliance numerators).
\bibitem{musk} N.~I. Muskhelishvili, \emph{Some Basic Problems of the
  Mathematical Theory of Elasticity} (point-force potentials).
\bibitem{rice} J.~R. Rice, ``Some Remarks on Elastic Crack-Tip Stress Fields,''
  Int.\ J.\ Solids Structures 8, 1972 (weight-function relations).
\bibitem{harter} J.~A. Harter, A.~V. Litvinov, ``Modeling Bearing Load in Wide
  Panels Using AFGROW,'' AFGROW European Workshop, 2016 (effective-width bearing).
\bibitem{newman} J.~C. Newman Jr., ``A crack opening stress equation for fatigue
  crack growth,'' Int.\ J.\ Fracture, 1984.
\end{thebibliography}

\end{document}
